\section{Introduction}

Ce laboratoire met en pratique les systèmes séquentiels synchrones vus en cours de SysLog1, et en particulier la décomposition en un bloc combinatoire d'état futur et un registre parallèle/parallèle. Des registres à décalage sont conçus sous Logisim, puis utilisés pour réaliser un chenillard sur 8 LEDs intégré dans la CPLD de la carte MAX\_V\_80P\_25P.

\begin{description}
	\item[Partie 1] Registre à décalage à droite de 4 bits avec enable, réalisé avec 4 bascules D.
	\item[Partie 2] Registre à décalage de 8 bits avec chargement parallèle et décalage à droite/gauche, basé sur un registre parallèle/parallèle.
	\item[Partie 3] Chenillard effectuant une rotation cyclique d'un pattern à droite ou à gauche, testé sur la carte MAX\_V\_80P\_25P.
\end{description}
